% !TeX root = ../main.tex
% Fonti: Lezione 11_12, pp. 1--6; richiami a Lezioni 5 e 9--10.
\chapter{Somme, formula di Grassmann e complementi}
\fonti{Lezioni 11 e 12}

\section{Somma di sottospazi}
\begin{definizione}
Per due sottospazi $U,W\subseteq V$, la loro \emph{somma} è
\[
U+W=\{u+w:u\in U,\ w\in W\}.
\]
È il più piccolo sottospazio che contiene $U\cup W$.
\end{definizione}
Infatti $0_V=0_U+0_W$ è nella somma e
$(u_1+w_1)+(u_2+w_2)=(u_1+u_2)+(w_1+w_2)$ vi appartiene;
lo stesso vale per il prodotto per scalari. La somma contiene
$U$ e $W$ perché si può aggiungere $0_V$; qualunque sottospazio
che li contenga contiene anche tutte le somme $u+w$. In generale
$U+W$ è più grande dell'unione: gli assi di $\R^2$ hanno somma
$\R^2$, mentre la loro unione contiene soltanto i due assi.

\section{La formula di Grassmann}
\begin{teorema}[formula di Grassmann]
Se $U$ e $W$ sono sottospazi di dimensione finita, allora
\[
\dim(U+W)=\dim U+\dim W-\dim(U\cap W).
\]
\end{teorema}
\begin{proof}
Sia $v_1,\ldots,v_r$ una base di $U\cap W$. Completiamola
rispettivamente a basi
\[
v_1,\ldots,v_r,u_1,\ldots,u_s\quad\text{di }U,
\qquad
v_1,\ldots,v_r,w_1,\ldots,w_t\quad\text{di }W.
\]
La famiglia unita
$v_1,\ldots,v_r,u_1,\ldots,u_s,w_1,\ldots,w_t$
genera $U+W$: ogni $u+w$ si ottiene espandendo $u$ e $w$ nelle
rispettive basi. Mostriamo che è indipendente. Supponiamo
\[
\sum_i\alpha_i v_i+\sum_j\beta_j u_j+
\sum_k\gamma_k w_k=0_V.
\]
Spostando l'ultima somma a destra, il membro sinistro appartiene
a $U$ e il destro a $W$, quindi il loro valore comune appartiene
a $U\cap W$. Esistono perciò $\delta_i$ con
$-\sum_k\gamma_k w_k=\sum_i\delta_i v_i$. L'indipendenza della
base di $W$ impone $\gamma_k=\delta_i=0$. La relazione iniziale
si riduce a $\sum_i\alpha_i v_i+\sum_j\beta_j u_j=0$, e la base
di $U$ impone anche $\alpha_i=\beta_j=0$.
La famiglia unita è dunque una base di $U+W$, che ha dimensione
$r+s+t$. Poiché $\dim U=r+s$, $\dim W=r+t$ e
$\dim(U\cap W)=r$, la formula segue.
\end{proof}

\section{Complementi e somma diretta}
\begin{definizione}
La somma $U+W$ si dice \emph{diretta}, e si scrive $U\oplus W$,
se $U\cap W=\{0_V\}$. Un \emph{complementare} di $U$ in $V$
è un sottospazio $W$ tale che $V=U\oplus W$.
\end{definizione}
La condizione di somma diretta equivale all'unicità della
scomposizione $v=u+w$. Se $u+w=u'+w'$, allora
$u-u'=w'-w\in U\cap W$; quando l'intersezione è nulla,
$u=u'$ e $w=w'$. Viceversa, un vettore non nullo nell'intersezione
avrebbe due scomposizioni, $x=x+0_V=0_V+x$.

In dimensione finita un complementare esiste sempre: si estende
una base di $U$ a una base di $V$ e si prende il sottospazio
generato dai vettori aggiunti. Esso non è in generale unico.
Per esempio, per l'asse orizzontale di $\R^2$ sono complementari
sia l'asse verticale sia ogni altra retta per l'origine distinta
dall'asse orizzontale.
\begin{schema}
\begin{tikzpicture}[scale=.7]
  \draw[asse] (-3,0)--(3.2,0) node[right] {$U$};
  \draw[spazio] (0,-2)--(0,2) node[above] {$W_1$};
  \draw[spazio] (-1.8,-1.8)--(1.8,1.8) node[above right] {$W_2$};
  \fill (0,0) circle (1.4pt);
\end{tikzpicture}
\end{schema}

\section{Un esercizio in \texorpdfstring{$\R^5$}{R5}}
\begin{esercizio}
Sia $V\subseteq\R^5$ l'iperpiano
$x_1+x_2+x_3+x_4+x_5=0$ e sia $U$ generato da
\begin{align*}
u_1&=(3,0,1,-2,-2),\\
u_2&=(1,-2,3,-1,-1),\\
u_3&=(3,6,-7,-1,-1).
\end{align*}
Si cercano una base e le equazioni cartesiane di $U$, un
complementare di $U$ in $V$ e un sottospazio $Z\subseteq\R^5$
di dimensione $3$ che intersechi $U$ e il complementare in
sottospazi di dimensione $1$.

Per prima cosa i tre generatori appartengono a $V$, perché la
somma delle coordinate di ciascuno è zero. Risolvendo
$\alpha_1u_1+\alpha_2u_2+\alpha_3u_3=0$, la seconda
coordinata dà $\alpha_2=3\alpha_3$ e la terza, insieme alla
prima, dà $\alpha_1=-2\alpha_3$. Quindi
\[
-2u_1+3u_2+u_3=0,
\qquad u_3=2u_1-3u_2.
\]
I vettori $u_1,u_2$ sono indipendenti: la seconda coordinata
di $a u_1+b u_2=0$ dà $b=0$ e la terza dà $a=0$.
Pertanto $(u_1,u_2)$ è una base di $U$ e $\dim U=2$.
La stessa relazione si può ottenere eliminando le righe della
matrice aumentata
\[
\left(\begin{array}{rrrrr|rrr}
3&0&1&-2&-2&1&0&0\\
1&-2&3&-1&-1&0&1&0\\
3&6&-7&-1&-1&0&0&1
\end{array}\right).
\]
Le operazioni $R_2\leftarrow R_1-3R_2$ e
$R_3\leftarrow R_3-R_1$ producono nella parte sinistra due
righe uguali $(0,6,-8,1,1)$; sottraendole, la riga nulla
porta nella parte destra i coefficienti $(-2,3,1)$.

Per le equazioni cartesiane scriviamo
$x=\lambda_1u_1+\lambda_2u_2$:
\[
\begin{array}{lll}
x_1=3\lambda_1+\lambda_2,&x_2=-2\lambda_2,&
x_3=\lambda_1+3\lambda_2,\\
x_4=-2\lambda_1-\lambda_2,&
x_5=-2\lambda_1-\lambda_2.&
\end{array}
\]
Da $\lambda_2=-x_2/2$ e
$\lambda_1=x_3+3x_2/2$ segue
\[
\boxed{\quad x_1-4x_2-3x_3=0,\qquad
2x_4+5x_2+4x_3=0,\qquad x_5-x_4=0.\quad}
\]
Queste tre equazioni sono anche sufficienti: dati $x_2,x_3$,
fissano $x_1,x_4,x_5$ esattamente come la parametrizzazione.

Per costruire un complementare in $V$ scegliamo
$w_1=(1,0,0,0,-1)$ e $w_2=(0,1,0,-1,0)$.
Entrambi sono in $V$. Verifichiamo l'indipendenza di
$u_1,u_2,w_1,w_2$: se
$a u_1+b u_2+c w_1+d w_2=0$, la terza coordinata dà
$a=-3b$, la seconda $d=2b$ e la quarta
$-2a-b-d=0$, cioè $3b=0$. Quindi $a=b=c=d=0$.
Poiché $V$ ha dimensione $4$ (una coordinata si ricava dalle
altre quattro), questi quattro vettori ne formano una base.
Dunque $W=\langle w_1,w_2\rangle$ soddisfa
$V=U\oplus W$.

Infine prendiamo
$Z=\langle u_1,w_1,e_1\rangle$, con
$e_1=(1,0,0,0,0)$. I tre vettori sono indipendenti:
in $a u_1+b w_1+c e_1=0$ la terza coordinata impone $a=0$,
la quinta $b=0$ e la prima $c=0$. Quindi $\dim Z=3$.
Ogni vettore di $Z$ ha seconda coordinata nulla. Se appartiene
a $U=\langle u_1,u_2\rangle$, il coefficiente di $u_2$ è
forzato a zero; pertanto $Z\cap U=\langle u_1\rangle$.
Se appartiene a $W=\langle w_1,w_2\rangle$, il coefficiente
di $w_2$ è zero; pertanto $Z\cap W=\langle w_1\rangle$.
In entrambi i casi la dimensione è $1$.
\end{esercizio}

\section{Equazioni di polinomi in basi diverse}
\begin{esercizio}
Sia $V=\R[x]_{\leq3}$ e
$W=\{p\in V:p(1)=p(2)=0\}$. Nella base monomiale
$(1,x,x^2,x^3)$, un polinomio
$p(x)=x_0+x_1x+x_2x^2+x_3x^3$ appartiene a $W$ se e
solo se le sue coordinate soddisfano
\[
x_0+x_1+x_2+x_3=0,\qquad
x_0+2x_1+4x_2+8x_3=0.
\]
Le due equazioni sono indipendenti, quindi $\dim W=2$.
Si può anche vedere direttamente che
$p(x)=(x-1)(x-2)(a+bx)$, con $a,b\in\R$.

Consideriamo ora la base
$(1,\,1-x,\,x-x^2,\,x^2-x^3)$ e scriviamo
\[
p(x)=y_0+y_1(1-x)+y_2(x-x^2)+y_3(x^2-x^3).
\]
È una base perché il passaggio alle coordinate monomiali ha
matrice triangolare con diagonale $(1,-1,-1,-1)$, quindi
invertibile. Valutando direttamente si ottiene
$p(1)=y_0$ e $p(2)=y_0-y_1-2y_2-4y_3$. Le equazioni
cartesiane di $W$ nelle nuove coordinate sono dunque
\[
y_0=0,\qquad y_0-y_1-2y_2-4y_3=0.
\]
L'insieme dei polinomi è lo stesso: cambiano le coordinate e
perciò cambia la forma delle equazioni.
\end{esercizio}
